\chapter{Aesthetics e Dynamics}
\label{g:aesthetics}
Estetica fantasy-dungeon sobria: la roadmap \`e la mappa del dungeon, i
topic sono stanze da conquistare, i boss sono i guardiani di capitolo con
nome e lore propri (Man-in-the-Middle, The Amnesiac, Spaghetti Colossus).
Le rarit\`a delle carte usano codice colore; le animazioni sono
volutamente sobrie per non distrarre dallo studio.

Dynamics: la progressione \`e doppia ma semplice -- \textbf{mastery}
(XP, livelli, badge, sblocco nodi) e \textbf{potere} (deck pi\`u forte
tramite ricompense 1-di-3), cos\`i che studiare meglio significhi
combattere meglio. La soglia 80\,\% mantiene il flow senza frustrare
(tentativi illimitati ma penalizzati); il boss a 10 HP contro 3 HP del
giocatore con quiz-danno a fine turno crea archi di tensione brevi e
leggibili. Nessun dato di esperienza utente \`e ancora raccolto: si veda
il Cap.~\ref{g:valutazione}.
